% !TEX root = ../main.tex
\chapter{Multilinear Algebra}\label{ch:c03}
\begin{paracol}{2}

\begin{sources}
\begin{tabularx}{\linewidth}{@{}l l L@{}}
\toprule
\textbf{Source} & \textbf{Length} & \textbf{Content}\\
\midrule
\seg{3}{1}{Lecture 3: Multilinear Algebra} & 1:42:36 & Why vector spaces although spacetime is not one; the vector space axioms and the abstract example of polynomials; ``what is a vector?'';
linear maps and their composition; the vector space $\operatorname{Hom}(V,W)$; the dual space and covectors; $(r,s)$-tensors as multilinear maps; vectors and covectors as tensors;
why not to choose a basis, and what a basis (Hamel basis) and the dimension are; components of vectors; the dual basis; components of tensors; the Einstein summation convention.\\
\seg{3}{2}{Tutorial 3: Multilinear Algebra} & 0:57:19 & Solutions of the sheet by a tutor (not Schuller): fifteen true/false statements; $\R^3$ as a vector space, a linear map and a covector on it,
the link to polynomials; components of tensors and (anti)symmetrisation of indices; linear maps as $(1,1)$-tensors.\\
\bottomrule
\end{tabularx}

\smallskip
\textbf{Transcripts} \file{materials/transcripts/C03-V01.txt} and \file{C03-V02.txt} (YouTube captions).
\textbf{Sheet 3:} see \url{https://tales.mbivert.com/on-heraeus-winter-school-tutorials/}.
\end{sources}

\begin{reading}[Where this chapter sits in the other notes]
\emph{Geometric Anatomy of Theoretical Physics} notes: ``Tensor space theory I'' (chapter 8: vector spaces, linear maps, the dual space, tensors as multilinear maps, bases,
components and the Einstein convention) and ``Tensor space theory II'' (chapter 11: the finite-dimensional case, the canonical isomorphisms and the tensor product treated
more slowly). The next lecture (chapter 4 here) goes on to differentiable manifolds, treated in ``Differentiable structures: definition and classification'' (chapter 7 there).
The tangent spaces of lecture 5, which are the vector spaces to which everything below will be applied, come in chapter 10 there.
\end{reading}

% ==================================================================
\section{Why vector spaces at all}\label{sec:c03-why}

Welcome back to the third lecture \ts{3}{1}{0:00}. Multilinear algebra is the same field as linear algebra: one just emphasises that, besides linear maps, one also studies \emph{multi}linear maps; the object of study
is vector spaces \ts{3}{1}{0:23}\ts{3}{1}{0:44}.

A word of warning from the beginning: \emph{spacetime will not be equipped with a vector space structure.} Before, we had a set and equipped it with a topology, and that is one step towards the picture of spacetime;
today we equip a set with a vector space structure, but this does not further the aim of describing spacetime. Schuller stresses it because it is often assumed that physical space or spacetime carries a vector space structure
\ts{3}{1}{1:08}\ts{3}{1}{1:29}\ts{3}{1}{1:49}.
\begin{aside}[Five times Paris]
``Where is five times Paris? Does that make any sense? Where does Paris plus Vienna lie?'' There is no notion on physical position space of adding places or of multiplying them by real numbers
\ts{3}{1}{2:10}\ts{3}{1}{2:32}. A student replies that one can define it, e.g.\ by multiplying each coordinate by five; Schuller: ``I tell you that at seven times Paris I hid a treasure of 1 million euros, it is really true,
go and find it: where do you go?'' The point is that one would first have to define the operations, and that one will find nothing in the world that justifies the definition; ``we can discuss this for half an hour and
I am sure in the end you will agree, but let me proceed; I beg you'' \ts{3}{1}{18:34}\ts{3}{1}{18:55}\ts{3}{1}{19:16}\ts{3}{1}{19:36}.
\end{aside}

So why study vector spaces? Because of the \emph{tangent spaces}. Their name is $T_pM$ (lecture 5, on Friday), they are constructed for smooth manifolds (lecture 4, today's next lecture). Very roughly (one will define them differently):
on a sphere one can attach a tangent space at the North Pole $p$, another at any other point $q$; many different tangent spaces. Whatever way one defines them, they carry a vector space structure, and a natural one at that:
the underlying space(time) is not a vector space, the tangent spaces are. And once one has a vector space, a derived notion arises, the notion of a \emph{tensor}, which is very important in differential geometry
\ts{3}{1}{2:32}\ts{3}{1}{2:53}\ts{3}{1}{3:14}\ts{3}{1}{3:35}\ts{3}{1}{3:56}\ts{3}{1}{4:17}.

It is best to study vector spaces and tensors first abstractly, for two reasons \ts{3}{1}{4:38}\ts{3}{1}{4:59}: (i) even the construction of $T_pM$, a finite-dimensional vector space, needs an intermediate vector space, called
$C^\infty(M)$ (``whatever that will be''), so one cannot construct it without knowing abstractly what vector spaces are; (ii) tensor techniques are most easily understood in an abstract setting, where one does not have to worry about a specific
vector space and its details \ts{3}{1}{5:20}\ts{3}{1}{5:42}\ts{3}{1}{6:02}\ts{3}{1}{6:23}. For a moment, therefore, we leave our aim where it is \ts{3}{1}{6:23}\ts{3}{1}{6:44}.

\section{Vector spaces}\label{sec:c03-vs}

Everyone will have seen the definition; nevertheless it is repeated \ts{3}{1}{6:44}.

\begin{definition}[New concepts: (real) vector space]\label{def:c03-vs}
A \textbf{vector space} is a triple $(V,+,\cdot)$ consisting of
\begin{enumerate}
  \item a set $V$ (not called $M$, since it is not the manifold $M$ from before; later $V$ will be, for most of the course, a tangent space $T_pM$ at one point or another);
  \item an \textbf{addition}, a map $+\colon V\times V\to V$ (it takes two vectors to a vector);
  \item an \textbf{s-multiplication}, a map $\cdot\colon\R\times V\to V$ (for a \emph{real} vector space; it takes a real number and a vector to a vector; it is not a multiplication between two vectors, which has nothing to do with the vector space structure.
        ``S'' may stand for scalar, but Schuller never uses that word)
\end{enumerate}
satisfying eight axioms. The first four concern the addition only (``K, M and Adu-Key'': commutative, associative, neutral element, inverse; the ``musketeer principle'' of the quantifiers ``for all'' and ``there exists''):
\begin{enumerate}
  \item[(a1)] commutativity: $v+w=w+v$ for all $v,w\in V$;
  \item[(a2)] associativity: $(u+v)+w=u+(v+w)$ for all $u,v,w\in V$;
  \item[(a3)] there exists a neutral element $\circ\in V$ (read: zero, but not the zero of $\R$) such that $v+\circ=v$ for all $v\in V$;
  \item[(a4)] for all $v\in V$ there exists a vector, written $-v$ (just a symbol, one only has to ensure that it lies in $V$), such that $v+(-v)=\circ$ (hence this law comes after the neutral element; $v-v=\circ$ is a short notation).
\end{enumerate}
The other four connect the two operations, with $\lambda,\mu\in\R$ and $v,w\in V$:
\begin{enumerate}
  \item[(s1)] $(\lambda\cdot\mu)\cdot v=\lambda\cdot(\mu\cdot v)$;
  \item[(s2)] $(\lambda+\mu)\cdot v=\lambda\cdot v+\mu\cdot v$;
  \item[(s3)] $\lambda\cdot(v+w)=\lambda\cdot v+\lambda\cdot w$;
  \item[(s4)] $1\cdot v=v$ (``a little trivial but important''; the $1$ is the one of $\R$)
\end{enumerate}
\ts{3}{1}{7:04}\ts{3}{1}{7:26}\ts{3}{1}{7:46}\ts{3}{1}{8:06}\ts{3}{1}{8:28}\ts{3}{1}{8:48}\ts{3}{1}{9:09}\ts{3}{1}{9:29}\ts{3}{1}{9:50}\ts{3}{1}{10:10}\ts{3}{1}{10:31}\ts{3}{1}{10:51}\ts{3}{1}{11:13}\ts{3}{1}{13:18}.
\end{definition}
In (s1) to (s3) one has to be careful about which plus and which times is which (it is worth marking them in red): in (s1) the first $\cdot$ on the left and the last $\cdot$ on the right are the s-multiplication on $V$, whereas $\lambda\cdot\mu$ is the multiplication of
real numbers; in (s2) the $+$ in $\lambda+\mu$ is the one of $\R$, the other is the one of $V$; in (s3) both $+$ are the one of $V$. One uses the same symbols in all cases, but at the beginning colours help
\ts{3}{1}{11:13}\ts{3}{1}{11:34}\ts{3}{1}{11:55}\ts{3}{1}{12:15}\ts{3}{1}{12:37}\ts{3}{1}{12:58}. The eight axioms are over a real number field here; everything can be slightly generalised to any underlying number field
\ts{3}{1}{36:36}.

\paragraph{``What is a vector?''} A word of terminology: an element of a vector space is informally called a \emph{vector}. One must not ask ``what is a vector?'': one can, and the answer is ``an element of a vector space'';
but if one hands over an object and asks whether it is a vector or not, one cannot say, because it depends on whether a vector space structure is defined in the background \ts{3}{1}{13:40}\ts{3}{1}{14:03}\ts{3}{1}{14:23}.

\subsection{The running example: polynomials}

Everyone has $\R^3$ in mind as the prime example of a vector space, but $\R^3$ is not a vector space all by itself: it can be \emph{made} into one. Citing Machiavelli (``The Prince'': one has to do the cruel things at the beginning),
Schuller will \emph{not} give an example that makes immediate contact with what one knows: one would imagine more than what he writes down. Instead, one example is used for everything constructed today: perfect, but not immediately
``oh, that is what I always thought vectors were''; this forces one to sharpen the understanding of the definitions. In the tutorials one is reconciled with the abstract notion, and $\R^n$ appears
\ts{3}{1}{14:23}\ts{3}{1}{14:44}\ts{3}{1}{15:04}\ts{3}{1}{15:24}\ts{3}{1}{15:45}.

\begin{example}[The set of polynomials $P$]\label{ex:c03-P}
Define the set
\[
  P:=\Bigl\{\,p\colon(-1,1)\to\R\ \Bigm|\ p(x)=\sum_{n=0}^{N}p_n\,x^n\ \text{with}\ p_n\in\R\Bigr\},
\]
the polynomials of degree at most a fixed $N$ (Schuller first says $N=7$, e.g.\ $x\mapsto x^2$ or $x\mapsto 2+x^7$ are in $P$ if $N\ge7$; later $N=3$ and, in the tutorial, $N=2$), defined on the real interval $(-1,1)$
\ts{3}{1}{15:45}\ts{3}{1}{16:06}\ts{3}{1}{16:27}\ts{3}{1}{16:48}\ts{3}{1}{17:08}.
\end{example}
At this stage a simple question: is $\mathrm{square}\colon x\mapsto x^2$ a vector? ``You should get the resounding answer: no.'' \texttt{square} is an element of the set $P$, but there is no addition and no s-multiplication defined anywhere around, so one
cannot even talk about a vector if one did not talk about a vector space in the first place \ts{3}{1}{17:08}\ts{3}{1}{17:31}\ts{3}{1}{17:52}.

One \emph{can} define an addition and an s-multiplication on $P$ (any way one wants, provided one can show the axioms afterwards): for $p,q\in P$ the sum $p\oplus q$ is the polynomial defined by
\[
  (p\oplus q)(x):=p(x)+q(x)
\]
(``you say: you wrote nothing new, the sum is the sum; but the left side is the sum on $P$ and the right-hand plus is the addition of real numbers, because $p(x)\in\R$''; the sum of polynomials is the polynomial whose value is the sum of the values, ``that is the trick''),
and for $\lambda\in\R$ the s-multiplication by $(\lambda\odot p)(x):=\lambda\cdot p(x)$ \ts{3}{1}{19:36}\ts{3}{1}{19:57}\ts{3}{1}{20:17}\ts{3}{1}{20:37}\ts{3}{1}{20:58}\ts{3}{1}{21:18}\ts{3}{1}{21:39}.
One checks the axioms (``and I am sure it is easy''): $(P,\oplus,\odot)$ is a vector space, and from now on the circles are dropped.

Now the question ``is \texttt{box} (some polynomial) a vector?'' arises again. One cannot look at it and say so: it is just an element of a set that may underlie a vector space; if someone insists that it is a vector but cannot tell how to add it to
another vector, ``he or she is in trouble''. Being an element of something only refers to the underlying set; ``it is a bad question whether something is a vector, but one can certainly talk about vector spaces''
\ts{3}{1}{21:39}\ts{3}{1}{22:00}\ts{3}{1}{22:21}\ts{3}{1}{22:42}\ts{3}{1}{23:02}\ts{3}{1}{23:23}.
In relativity the same question is asked of tensors: ``what is a tensor?'' How does one see that something is a tensor? One does not: one can only define \emph{tensor spaces}, and informally call their elements tensors (vectors are an example).
Schuller would rather discourage one from using the word at all formally, ``of course I myself will not adhere to this advice'' \ts{3}{1}{23:23}\ts{3}{1}{23:43}.

\paragraph{Student question: a quantum state is a vector in a Hilbert space?} How would Schuller rephrase ``a quantum mechanical state is a vector in a Hilbert space''? He would say it is an \emph{element} of the Hilbert space; and
strictly a state is a ray in the Hilbert space, not a vector, because of the ambiguity of phase and norm. Why call it a vector? \ts{3}{1}{24:06}\ts{3}{1}{24:26}\ts{3}{1}{24:46}.

\section{Linear maps}\label{sec:c03-linear}

A standard procedure in mathematics, once a structure is defined, is to study it by maps between two instances of it that respect the structure; in topology these were the continuous maps (pre-images of open sets are open). For vector spaces they are the
\emph{linear maps}; ``this is always a good idea to learn more at any structural level'' \ts{3}{1}{25:08}\ts{3}{1}{25:28}\ts{3}{1}{25:49}\ts{3}{1}{26:10}.

\begin{definition}[New concepts: linear map]\label{def:c03-linear}
Let $(V,+_V,\cdot_V)$ and $(W,+_W,\cdot_W)$ be vector spaces (any two; the operations need have nothing to do with each other). A map $f\colon V\to W$ is called \textbf{linear} if
\begin{enumerate}
  \item $f(v+_V\tilde v)=f(v)+_Wf(\tilde v)$ for all $v,\tilde v\in V$;
  \item $f(\lambda\cdot_V v)=\lambda\cdot_W f(v)$ for all $v\in V$, $\lambda\in\R$.
\end{enumerate}
One writes $f\colon V\xrightarrow{\ \sim\ }W$ (a tilde on the arrow) to say ``linear map between vector spaces'' without having to specify the operations each time
\ts{3}{1}{26:31}\ts{3}{1}{26:52}\ts{3}{1}{27:12}\ts{3}{1}{33:08}.
\end{definition}
All four operations enter the definition: the plus on $V$, the plus on $W$, the s-multiplications on $V$ and on $W$. As in topology, where one can speak of continuity of a map only once the sets carry a topology, one can speak of linearity only
once there are vector space structures on both sets: given a map, one cannot ask whether it is linear unless one is told the addition and s-multiplication on domain and target \ts{3}{1}{27:34}\ts{3}{1}{27:55}\ts{3}{1}{28:16}\ts{3}{1}{28:37}.
This holds also for $\R^3$: one must say how to add two triples of reals (entrywise) and how to multiply by a scalar (pull it into each entry); only then has $\R^3$ been made into a vector space \ts{3}{1}{28:37}\ts{3}{1}{28:57}\ts{3}{1}{29:18}.
\begin{coursenote}[Nothing natural about entrywise addition]
As with the standard topology, one is ``indoctrinated'' to add triples componentwise, so that it seems a natural choice; but \emph{canonical} in mathematics means: uniquely constructible from the structure already given. If one has to say more (``take it
componentwise'', ``pull the $\lambda$ into each entry''), that is additional information, and additional information is structure. ``There is nothing natural about adding two elements of $\R^3$ by adding the entries; you may call it natural because you are used to it, but that
depends on your history and that is not a good definition of natural'' \ts{3}{1}{29:18}\ts{3}{1}{29:38}\ts{3}{1}{30:00}\ts{3}{1}{30:21}.
\end{coursenote}

\begin{example}[Differentiation on $P$]\label{ex:c03-delta}
Define $\Delta\colon P\xrightarrow{\ \sim\ }P$ by $\Delta(p):=p'$, the derivative of $p$ (a polynomial of degree at most $N$ again). It is linear: $\Delta(p+q)=(p+q)'=p'+q'=\Delta(p)+\Delta(q)$ (the sum rule, and the plus is the one of $P$), and $\Delta(\lambda p)=\lambda p'=\lambda\Delta(p)$
\ts{3}{1}{30:41}\ts{3}{1}{31:02}\ts{3}{1}{31:23}\ts{3}{1}{31:44}\ts{3}{1}{32:05}\ts{3}{1}{32:26}\ts{3}{1}{32:47}.
\end{example}

\paragraph{Composition.} Let $f\colon V\to W$ and $g\colon W\to U$ be linear. Then $g\circ f\colon V\to U$ is linear: ``the diagram says it all, you can go this way or that way'' \ts{3}{1}{33:28}\ts{3}{1}{33:48}\ts{3}{1}{34:09}.
(If one is ``so arbitrary as to introduce a basis'', this fact explains the matrix multiplication rule, once one decides how to store the information of a linear map; ``lots of the rules you know are very arcane rules with lots of information for very simple facts; this is the very simple fact
that explains everything about matrix multiplication'', but a basis is introduced only at the end of this lecture, because it ``clouds your thinking'' \ts{3}{1}{34:31}\ts{3}{1}{34:52}\ts{3}{1}{35:12}.)
In the example, $\Delta\circ\Delta\colon P\to P$ is the second derivative operator, still linear, and so are all the higher ones \ts{3}{1}{35:12}\ts{3}{1}{35:32}.

\subsection{The vector space of linear maps}

``Fun fact'' (very useful) \ts{3}{1}{35:55}\ts{3}{1}{36:16}: let $V,W$ be vector spaces. One can collect all linear maps between them in a set, called
\[
  \operatorname{Hom}(V,W):=\{f\colon V\xrightarrow{\ \sim\ }W\}
\]
(from \emph{homomorphism}: a structure-preserving map; a funny name for a very simple set, the collection of all linear maps; the order of $V$ and $W$ matters) \ts{3}{1}{36:36}\ts{3}{1}{36:57}.
Is it a vector space? Not yet, but it can be made into one \ts{3}{1}{37:17}\ts{3}{1}{37:39}.

\begin{definition}[New concepts: the vector space $\operatorname{Hom}(V,W)$]\label{def:c03-hom}
For $f,g\in\operatorname{Hom}(V,W)$ define the sum $f\oplus g$ (a third plus, hence the circle; the circle is dropped later) pointwise, using the plus on $W$:
\[
  (f\oplus g)(v):=f(v)+_Wg(v)\qquad(v\in V),
\]
and similarly the s-multiplication $(\lambda\odot f)(v):=\lambda\cdot_Wf(v)$. Then $\operatorname{Hom}(V,W)$ with these operations is again a vector space (one checks the axioms, which is very simple)
\ts{3}{1}{37:39}\ts{3}{1}{38:00}\ts{3}{1}{38:20}\ts{3}{1}{38:41}\ts{3}{1}{39:02}.
\end{definition}
So from two vector spaces one immediately constructs a third. This could be a curiosity, but it is a \emph{tower of structures}: $\operatorname{Hom}(V,W)$ is a vector space, so one can look at linear maps from it into another vector space,
$\operatorname{Hom}(\operatorname{Hom}(V,W),U)$, and so on; the vector space structure inherits itself again and again; ``we will come back to this'' \ts{3}{1}{39:22}\ts{3}{1}{39:43}.

\begin{example}[$\operatorname{Hom}(P,P)$: mixed derivative operators]\label{ex:c03-homPP}
$\operatorname{Hom}(P,P)$ is a vector space; $\Delta$ and $\Delta\circ\Delta$ and every $m$-th derivative are elements of it. Since it is a vector space one can form, with the plus and the s-multiplication of $\operatorname{Hom}(P,P)$,
for instance $5\cdot\Delta+\Delta\circ\Delta$, ``five times the first derivative plus the second derivative'', and one does not need to think: the result lies in $\operatorname{Hom}(P,P)$ again. If one is ever in doubt whether mixed-order derivative operators are still
linear, they are; ``this is something you do all the time'' \ts{3}{1}{40:04}\ts{3}{1}{40:25}\ts{3}{1}{40:45}\ts{3}{1}{41:05}\ts{3}{1}{41:26}\ts{3}{1}{41:47}.
\end{example}

\section{The dual vector space}\label{sec:c03-dual}

``Now comes a much feared topic, the dual vector space, and it is totally unclear why it is so feared, because we already did it: it is a heavily used special case of the stuff before'' \ts{3}{1}{42:08}\ts{3}{1}{42:29}.

\begin{definition}[New concepts: dual vector space, covector]\label{def:c03-dual}
For a vector space $V$ the \textbf{dual vector space} (the dual of $V$) is
\[
  V^*:=\operatorname{Hom}(V,\R)=\{\varphi\colon V\xrightarrow{\ \sim\ }\R\},
\]
the set of linear maps from $V$ into the reals, with the plus and s-multiplication defined as above \ts{3}{1}{42:29}\ts{3}{1}{42:50}\ts{3}{1}{43:11}\ts{3}{1}{44:13}\ts{3}{1}{44:34}.
An element $\varphi\in V^*$ is informally called a \textbf{covector} (``co-vector'') \ts{3}{1}{44:55}\ts{3}{1}{45:16}.
\end{definition}
Here $\R$ with its addition and multiplication is a real vector space (one quickly checks the axioms), so one can speak of linearity. It is just a very special case of $\operatorname{Hom}(V,W)$, but heavily used, because if one knows $V$ and $V^*$ one can construct all the rest:
$V$ and $V^*$ are the elementary building blocks, and ``it is always clever to break something down to its most elementary building blocks'' \ts{3}{1}{42:50}\ts{3}{1}{43:31}\ts{3}{1}{43:52}.
(``You say, it is a very abstract set, a set of maps: well, it is not so bad; you understand the set of polynomials, that is even more complicated.'')

\paragraph{Terminology.} It is as bad as for ``vector''. $V^*$ is a vector space in its own right, so a covector is also a vector, being an element of a vector space; the convention is a tacit agreement: given a vector space $V$, its elements are called vectors, the elements of the
corresponding $V^*$ covectors. The ``co'' only makes sense in relation to $V$; one cannot look at something and say ``this is a covector, not a vector''. Naturally one will be given a vector space, the tangent spaces to the manifold, whose elements are vectors, and whose dual elements are covectors
\ts{3}{1}{44:55}\ts{3}{1}{45:16}\ts{3}{1}{45:37}\ts{3}{1}{45:57}\ts{3}{1}{46:18}.

\begin{example}[Integration is a covector on $P$]\label{ex:c03-int}
Define $I\colon P\xrightarrow{\ \sim\ }\R$ by
\[
  I(p):=\int_0^1\mathrm dx\,p(x)
\]
(the interval is open, but ``it does not matter whether the end point is in there or not''). It is linear, by the rules of integration: $I(p+q)=I(p)+I(q)$ (split the integral) and $I(\lambda p)=\lambda I(p)$. So $I\in P^*$: the operation of definite integration is a covector.
In quantum mechanics one comes with differentiation operators from the left, or with an integral operator; ``it is nothing, a simple thing, it is a covector'' \ts{3}{1}{46:18}\ts{3}{1}{46:38}\ts{3}{1}{46:58}\ts{3}{1}{47:19}\ts{3}{1}{47:39}\ts{3}{1}{48:00}\ts{3}{1}{48:21}\ts{3}{1}{48:42}.
\end{example}
\begin{aside}[The gradient is a covector]
One has been indoctrinated that the gradient of a function on $\R^n$ is a vector. No: if vectors are directions, the gradient is a covector, and ``everything else gets you into trouble''; it can be seen with the multi-dimensional chain rule:
``if you hold that the gradient is a vector and you hold that the chain rule applies, you can prove anything, because it is inconsistent''. One has met lots of objects that are covectors but that were called vectors ``in order not to bother you with these details'';
the impression that covectors are unnatural or esoteric ``is simply because people have lied to you''; they are all around and one knows many examples without knowing. Several points are taken up again later
\ts{3}{1}{49:03}\ts{3}{1}{49:24}\ts{3}{1}{49:44}\ts{3}{1}{50:05}.
\end{aside}

\section{Tensors}\label{sec:c03-tensors}

Given a vector space, its dual is a very special case of maps from it into another vector space, and with these building blocks one has what is needed to build the whole structure announced above, in which one heaps maps on maps on maps and always
equips them with a vector space structure; the elements of these higher and higher structures are generically called \emph{tensors} \ts{3}{1}{50:32}\ts{3}{1}{50:52}. They are a natural object of linear algebra, ``you just need to write them down''.
Everything said so far holds for any vector space (any dimension, not defined yet), and later a dimension and a basis will be introduced, which however brings a great deal of arbitrariness, the choice of a basis. On finite-dimensional vector spaces there is a very
simple way to think of tensors, as \emph{multilinear maps}, and this is how they are introduced: the tangent spaces to an $d$-dimensional manifold will be finite dimensional, with vector space dimension $d$ (a different notion of dimension from the one of a topological manifold defined before, but
the two numbers agree); mostly, especially for general relativity, one will deal with four-dimensional vector spaces \ts{3}{1}{51:13}\ts{3}{1}{51:33}\ts{3}{1}{51:54}\ts{3}{1}{52:14}\ts{3}{1}{52:35}.

\begin{definition}[New concepts: $(r,s)$-tensor]\label{def:c03-tensor}
Let $(V,+,\cdot)$ be a vector space and $r,s\in\mathbb N_0$. An \textbf{$(r,s)$-tensor over $V$} is a multilinear map
\[
  T\colon\underbrace{V^*\times\dots\times V^*}_{r\ \text{copies}}\times\underbrace{V\times\dots\times V}_{s\ \text{copies}}\xrightarrow{\ \sim\ }\R,
\]
i.e.\ it eats $r$ covectors and $s$ vectors (the elements of the Cartesian product are tuples whose first $r$ entries are covectors and whose last $s$ entries are vectors) and returns a real number, linearly in each slot
\ts{3}{1}{52:58}\ts{3}{1}{53:21}\ts{3}{1}{53:41}\ts{3}{1}{54:03}\ts{3}{1}{54:24}.
\end{definition}
One is tempted to write $r+s$ tildes on the arrow; one writes just one. Multilinearity is best explained by an example: let $T$ be a $(1,1)$-tensor, so it eats a covector $\varphi\in V^*$ and a vector $v\in V$. Linearity in the first slot means
\[
  T(\varphi+\psi,v)=T(\varphi,v)+T(\psi,v),\qquad T(\lambda\varphi,v)=\lambda T(\varphi,v),
\]
and independently also in the second slot,
\[
  T(\varphi,v+w)=T(\varphi,v)+T(\varphi,w),\qquad T(\varphi,\lambda v)=\lambda T(\varphi,v).
\]
Using linearity twice,
\[
  T(\varphi+\psi,v+w)=T(\varphi,v)+T(\varphi,w)+T(\psi,v)+T(\psi,w)
\]
\ts{3}{1}{54:45}\ts{3}{1}{55:06}\ts{3}{1}{55:27}\ts{3}{1}{55:48}\ts{3}{1}{56:08}\ts{3}{1}{56:29}\ts{3}{1}{56:50}.

\paragraph{Student question: a $(1,1)$-tensor as a linear map $V\to V$.} A student associates with a $(1,1)$-tensor something one feeds a vector and gets a vector out. Schuller: it is precisely the same data (the question was reserved for the problem sheet, but it is a very good one). Excursion
\ts{3}{1}{57:12}\ts{3}{1}{57:34}:
given a multilinear $T\colon V^*\times V\to\R$, define a map $\phi_T\colon V\to V^{**}$ as follows: for a vector $v\in V$ put $\phi_T(v):=T(\,\cdot\,,v)$, the map that eats a covector $\varphi$ and returns the real number $T(\varphi,v)$. Once $v$ is fixed
this is linear in the open slot, so $\phi_T(v)$ is a linear map from $V^*$ to $\R$, i.e.\ an element of $V^{**}=(V^*)^*$, \emph{not} of $V$.
Now there is a fact, not proved here: if $\dim V<\infty$, then $V^{**}$ reproduces $V$ (the identification $v\mapsto(\varphi\mapsto\varphi(v))$, an isomorphism only for finite dimension). So after all a map $V\to V$ has been constructed, which is the way one likes to think of a $(1,1)$-tensor.
Otherwise one would have to put $r$ copies of $V^*$ and $s$ of $V$, then $p$ of $V^{**}$ and $q$ of $V^{***}$, and classify tensors by $(r,s,p,q)$; ``whenever I have a $V$ in a slot I can put it on the other side by adding a star, or a star can be removed, it always contains the same data''.
Conversely, given a linear map $\phi\colon V\to V$, construct $T_\phi\colon V^*\times V\to\R$, $T_\phi(\varphi,v):=\varphi(\phi(v))$: take the covector and the vector, map $v$ by $\phi$ and apply $\varphi$.
From this one can show that $T$ is recovered from the $\phi_T$ constructed from it, and $\phi$ from the $T_\phi$ constructed from $\phi$: for finite-dimensional $V$ the two descriptions contain the same data, ``but this is a canonical notation for it''
\ts{3}{1}{57:34}\ts{3}{1}{57:54}\ts{3}{1}{58:16}\ts{3}{1}{58:36}\ts{3}{1}{58:56}\ts{3}{1}{59:17}\ts{3}{1}{59:37}\ts{3}{1}{59:59}\ts{3}{1}{1:00:19}\ts{3}{1}{1:00:40}\ts{3}{1}{1:01:01}\ts{3}{1}{1:01:22}\ts{3}{1}{1:01:42}\ts{3}{1}{1:02:03}\ts{3}{1}{1:02:23}.
(The first claim above is Schuller's, the explicit formula $T_\phi(\varphi,v)=\varphi(\phi(v))$ is the natural reading of his words [?]; the whole construction is Exercise~4 of the tutorial, \cref{sec:t03-1}.)

\begin{example}[The inner product on $P$ is a $(0,2)$-tensor]\label{ex:c03-ip}
Define $g\colon P\times P\xrightarrow{\ \sim\ }\R$ by
\[
  g(p,q):=\int_{-1}^{1}\mathrm dx\,p(x)\,q(x).
\]
It is multilinear, as one splits the integral (``so this is a $(0,2)$ tensor over $P$''): it takes two polynomials and no covectors. It is the inner product on function space as it appears in quantum mechanics (where one needs more functions than polynomials and, being complex valued, a star on the first slot).
A linear map $P\to P$ as before is a $(1,1)$-tensor \ts{3}{1}{1:02:46}\ts{3}{1}{1:03:06}\ts{3}{1}{1:03:27}\ts{3}{1}{1:03:48}\ts{3}{1}{1:04:08}\ts{3}{1}{1:04:29}\ts{3}{1}{1:04:49}.
\end{example}
Schuller promises that on finite-dimensional vector spaces at least, all the structures that heap on each other and inherit their vector space structure can be constructed from the elementary building blocks $V$ and $V^*$, ``and that is why they are so important''
\ts{3}{1}{1:04:49}\ts{3}{1}{1:05:10}. (A remark about $T\in\operatorname{Hom}(V,W)$ is cut off: ``it can be done but it leads too far'' \ts{3}{1}{1:05:10}\ts{3}{1}{1:05:31}.)

\subsection{Vectors and covectors are tensors}

\begin{theorem}[Covectors are $(0,1)$-tensors; vectors are $(1,0)$-tensors]\label{thm:c03-vec-tensors}
\begin{enumerate}
  \item A covector $\varphi\in V^*$ is a $(0,1)$-tensor.
  \item If $\dim V<\infty$, a vector $v\in V$ can be considered a $(1,0)$-tensor.
\end{enumerate}
\end{theorem}
\begin{proof}
(i) By definition of $V^*$, $\varphi$ is a linear map from $V$ to $\R$; by the definition of tensors this is a $(0,1)$-tensor (zero copies of $V^*$, one of $V$): ``theorem and proof in one''.
(ii) Recall (without proof) that for $\dim V<\infty$ one has $V=V^{**}$: from a given $v$ one constructs the linear map $V^*\to\R$, $\varphi\mapsto\varphi(v)$, which contains the same information, and from this map one reconstructs $v$ (a student remarks that one wants the isomorphism, and
that it would be enough to have an inclusion in one direction to \emph{say} this, plus the other direction to reconstruct; Schuller agrees ``so we agree it is true what I write, but I could write something weaker''). This linear map $V^*\to\R$ is by definition a $(1,0)$-tensor
\ts{3}{1}{1:05:53}\ts{3}{1}{1:06:15}\ts{3}{1}{1:06:38}\ts{3}{1}{1:07:01}\ts{3}{1}{1:07:21}\ts{3}{1}{1:07:42}\ts{3}{1}{1:08:03}\ts{3}{1}{1:08:23}.
\end{proof}
So vectors can be considered $(1,0)$-tensors, covectors $(0,1)$-tensors, linear maps $(1,1)$-tensors, and an inner product on a vector space (additional data, if one is given one) is a $(0,2)$-tensor. An \emph{inverse} of an inner product, ``it should be called the dual inner product'', would be a
$(2,0)$-tensor, an inner product on the dual space. Many examples have been seen in physics classes and probably never called tensors \ts{3}{1}{1:08:23}\ts{3}{1}{1:08:46}\ts{3}{1}{1:09:08}\ts{3}{1}{1:09:29}.

\begin{exercise}[The inertia tensor]\label{ex:c03-inertia}
Homework: the inertia tensor of classical mechanics, with two indices. Is it a $(1,1)$-, a $(2,0)$- or a $(0,2)$-tensor? ``Go to your mechanics book and look at it'' \ts{3}{1}{1:09:29}\ts{3}{1}{1:09:49}.
\end{exercise}
\begin{coursenote}[On the inertia tensor]
The kinetic energy of a rigid body, $2T=I(\omega,\omega)$, eats two angular velocities (vectors) and returns a number, so with that reading $I$ is a $(0,2)$-tensor; the same data viewed as a map $\omega\mapsto L$ is a map from $V$ into $V^*$ (angular momentum is naturally a covector).
(Not said in class.)
\end{coursenote}

\section{Bases: why one should avoid choosing one}\label{sec:c03-basis}

All of this has been said without ever talking about a basis, about the components of a vector, even without thinking of a vector as a collection of numbers (``in terms of the polynomials, of course, the numbers are hidden somewhere, but that is not the point; the point is not to see them''):
``if you do not see the numbers in the vectors, you did well, because that is the structure of the whole thing. Anything that can be constructed this way is natural, and you can make your proofs'' \ts{3}{1}{1:09:49}\ts{3}{1}{1:10:10}\ts{3}{1}{1:10:30}\ts{3}{1}{1:10:51}.

There is a notion of basis of a vector space, and special bases of the dual space (dual bases), and it has its place; but to associate to a vector a collection of numbers, the vector space must be finite dimensional and one must \emph{choose}. Given the polynomial $x^2+x^3$, which
numbers correspond to it? There are uncountably many ways to associate this one vector with a set of numbers, and a choice of one of these ways is called a choice of basis; ``it is a choice that has nothing to do with the object you wrote down'', as in quantum mechanics, where the inner product is what counts
\ts{3}{1}{1:10:51}\ts{3}{1}{1:11:11}\ts{3}{1}{1:11:31}\ts{3}{1}{1:11:52}. A choice out of uncountably many is a massive amount of information: one can make it and earn the fruits, as long as one keeps in mind that everything written down with it depends on this choice and has
nothing to do with the vector alone but with the vector and the choice. ``It is a very mixed blessing to choose a basis; nevertheless we will do it because it is quite practical for explicit calculations''; in England, ``a gentleman, and of course today a lady, only chooses a basis if he or she has to''
\ts{3}{1}{1:11:52}\ts{3}{1}{1:12:13}\ts{3}{1}{1:12:33}\ts{3}{1}{1:12:54}\ts{3}{1}{1:13:14}.

\begin{definition}[New concepts: (Hamel) basis]\label{def:c03-basis}
Let $V$ be a vector space. A subset $B\subseteq V$ is called a \textbf{basis} (a \emph{Hamel basis}, the basis of linear algebra) if for every $v\in V$ there exist a \emph{finite} subset $\{f_1,\dots,f_N\}\subseteq B$ and unique real numbers $v^1,\dots,v^N$ such that
\[
  v=v^1f_1+\dots+v^Nf_N.
\]
\ts{3}{1}{1:13:35}\ts{3}{1}{1:13:57}\ts{3}{1}{1:14:17}\ts{3}{1}{1:15:00}\ts{3}{1}{1:15:21}\ts{3}{1}{1:15:42}\ts{3}{1}{1:16:02}\ts{3}{1}{1:16:23}.
\end{definition}
($N$ may vary with $v$; $B$ itself may have infinitely many elements.)
Comments, as given in class:
\begin{itemize}
  \item In quantum mechanics bases often have infinitely many elements and one combines a vector by an infinite series; that is \emph{not} the basis of linear algebra, although most books tacitly assume it or do not mention the opposite
        \ts{3}{1}{1:14:17}\ts{3}{1}{1:14:39}. Hence the name Hamel basis for the linear algebra one. Series only converge if one has a notion of convergence, so any more liberal notion needs a topology on the vector space on top of the vector space structure; that is called a \emph{Schauder basis}
        \ts{3}{1}{1:18:16}\ts{3}{1}{1:18:37}\ts{3}{1}{1:18:57}.
  \item \emph{Student question on the harmonic oscillator.} The infinitely many energy eigenfunctions of the harmonic oscillator are not a Hamel basis of the Hilbert space but a Schauder basis, for which one needs additional structure: the Hilbert space has an inner product, a norm, and is complete, so the series converge.
        ``That is not the basis of linear algebra; it is a basis on a topological vector space, not on a bare vector space: it is a matter of how much structure exactly you assume'' \ts{3}{1}{1:18:57}\ts{3}{1}{1:19:18}\ts{3}{1}{1:19:39}.
  \item \emph{Student question on uniqueness.} Is there an exclamation mark (``exists a unique'') on the finite subset as well? Schuller: the numbers are unique; if another stretch of $B$ gave the same vector in another way, that would go against the linear independence of the basis; the way to be safe is to require that all
        coefficients are non-zero (the finite subset is then also unique), or to say that almost all coefficients are zero and a series becomes a sum. ``Let us keep it like this'' \ts{3}{1}{1:16:23}\ts{3}{1}{1:16:54}\ts{3}{1}{1:17:14}\ts{3}{1}{1:17:35}\ts{3}{1}{1:17:55}.
\end{itemize}

\begin{definition}[New concepts: dimension]\label{def:c03-dim}
If there exists a basis $B$ with finitely many elements, $d$ say, one calls $d=:\dim V$ the \textbf{dimension} of the vector space, and $V$ \textbf{finite dimensional}
\ts{3}{1}{1:19:39}\ts{3}{1}{1:20:00}\ts{3}{1}{1:20:22}.
\end{definition}
There is not a unique basis (there are uncountably many), so one has to show the definition is well defined: every other basis has the same number of elements (asserted here without proof). Once the notion is defined, one knows what a finite-dimensional vector space is; if one started from ``a finite basis''
one would never know whether one got them all. Infinite-dimensional vector spaces are used occasionally, but one will never choose a basis for them; for the finite-dimensional ones one will often do so \ts{3}{1}{1:20:22}\ts{3}{1}{1:20:42}\ts{3}{1}{1:21:03}\ts{3}{1}{1:21:23}\ts{3}{1}{1:21:44}.

\begin{remark}[Components of a vector]\label{rem:c03-components}
From now on (for simplicity) $V$ is finite dimensional. Having chosen a basis $e_1,\dots,e_d$ of $V$, one uniquely associates with every $v\in V$ a tuple $(v^1,\dots,v^d)\in\R^d$, the \textbf{components of $v$ with respect to the chosen basis}, namely the unique numbers such that
\[
  v=v^1e_1+\dots+v^de_d.
\]
A vector does not have components all by itself, only once a basis is chosen; ``you pay a big cost: everything will depend on an arbitrary choice'' \ts{3}{1}{1:21:44}\ts{3}{1}{1:22:04}\ts{3}{1}{1:22:26}\ts{3}{1}{1:22:48}\ts{3}{1}{1:23:09}\ts{3}{1}{1:23:29}\ts{3}{1}{1:23:50}\ts{3}{1}{1:24:12}.
\end{remark}

\paragraph{Student question: charts were arbitrary too, so why is the basis bad?} Introducing a chart (last lectures) is also an arbitrariness that depends on the publisher; why is it bad here and not there? It is bad there too; ``hence I would only buy an atlas from a publisher who sells a \emph{maximal} atlas with all the charts in it, and then make sure everything I define is independent of
my fantasy''. It is the same here: one can introduce a basis and define concepts using the objects defined in it, but one then has to check that under a change of basis the objects do not change. Example: the determinant of a linear map can be defined by choosing a basis and associating a matrix, and then one has to check invariance;
``my advice would be never to introduce a basis in the first place and define the determinant basis free''; the same for the trace. That is exactly the manifold philosophy of chapter~\ref{ch:c02} \ts{3}{1}{1:24:12}\ts{3}{1}{1:24:33}\ts{3}{1}{1:24:54}\ts{3}{1}{1:25:14}\ts{3}{1}{1:25:35}.

\subsection{The dual basis}

Having chosen a basis $e_1,\dots,e_d$ of $V$, one can also choose a basis $\epsilon^1,\dots,\epsilon^d$ (indices up) of $V^*$, entirely independently, $V^*$ being a different vector space; but this makes double the arbitrariness. It is more economical to require that the basis of $V^*$ be tied to the one of $V$
\ts{3}{1}{1:25:56}\ts{3}{1}{1:26:17}\ts{3}{1}{1:26:38}\ts{3}{1}{1:26:58}\ts{3}{1}{1:27:19}.

\begin{definition}[New concepts: dual basis]\label{def:c03-dualbasis}
Given a basis $e_1,\dots,e_d$ of $V$, a basis $\epsilon^1,\dots,\epsilon^d$ of $V^*$ is called the \textbf{dual basis} if
\[
  \epsilon^a(e_b)=\delta^a{}_b=\begin{cases}1&a=b\\0&a\neq b\end{cases}\qquad(a,b=1,\dots,d).
\]
This requirement determines the $\epsilon^a$ uniquely (a linear system, solvable in exactly one way), hence ``the'' dual basis; the whole purpose is to keep the choices down by linking the two; it is a matter of economy and, as an additional benefit, it makes the formulas simpler in many places
\ts{3}{1}{1:27:19}\ts{3}{1}{1:27:39}\ts{3}{1}{1:28:00}\ts{3}{1}{1:28:21}\ts{3}{1}{1:28:41}\ts{3}{1}{1:29:02}.
\end{definition}

\begin{example}[Dual basis of $P$ for $N=3$]\label{ex:c03-dualbasis}
Take $N=3$, so $\dim P=4$. A basis: $e_a(x):=x^a$ for $a=0,1,2,3$, i.e.\ the polynomials $1,x,x^2,x^3$ (every polynomial up to degree three is a linear combination of them; ``I could choose a totally different one''). The dual basis consists of the covectors
\[
  \epsilon^a(p):=\frac1{a!}\,\frac{\mathrm d^ap}{\mathrm dx^a}\Bigr|_{x=0},
\]
(``$0!=1$''): the $a$-th derivative divided by $a!$, evaluated at zero, a polynomial sent to a number. Indeed the elements of the dual basis are not polynomials, they are derivative operators with an evaluation. Check: $\epsilon^a(e_b)=\frac1{a!}\frac{\mathrm d^a}{\mathrm dx^a}x^b|_0$; the $a$-th derivative gives a factor $a!$ for $b=a$; for $b<a$ the derivative has already killed $x^b$;
for $b>a$ a power of $x$ remains which vanishes at $0$; so $\epsilon^a(e_b)=\delta^a{}_b$ \ts{3}{1}{1:29:02}\ts{3}{1}{1:29:26}\ts{3}{1}{1:29:46}\ts{3}{1}{1:30:07}\ts{3}{1}{1:30:28}\ts{3}{1}{1:30:50}\ts{3}{1}{1:31:10}\ts{3}{1}{1:31:31}\ts{3}{1}{1:31:52}\ts{3}{1}{1:32:12}\ts{3}{1}{1:32:33}\ts{3}{1}{1:32:54}.
(In the lecture the index $b$ is called ``the $b$-th order'' in the check and the case $b<a$ is the lower orders ``killed''.)
\end{example}
\begin{coursenote}[Taylor]
This is nothing but Taylor's formula: $p(x)=\sum_a\epsilon^a(p)\,x^a$, the components of $p$ with respect to the basis $(x^a)$ are the $\epsilon^a(p)$ (see \cref{rem:c03-comp-tensor}). (Not said in class.)
\end{coursenote}

\subsection{Components of tensors}

``Now comes something everybody has been waiting for who knew tensors before'': components of tensors \ts{3}{1}{1:32:54}\ts{3}{1}{1:33:15}.

\begin{definition}[New concepts: components of a tensor]\label{def:c03-comp-tensor}
Let $T$ be an $(r,s)$-tensor over the finite-dimensional vector space $V$, $e_1,\dots,e_d$ a basis of $V$ and $\epsilon^1,\dots,\epsilon^d$ its dual basis (``a tacit agreement; one can do it differently but one gets more work''). The $d^{\,r+s}$ real numbers
\[
  T^{i_1\dots i_r}{}_{j_1\dots j_s}:=T\bigl(\epsilon^{i_1},\dots,\epsilon^{i_r},e_{j_1},\dots,e_{j_s}\bigr),\qquad i_1,\dots,i_r,\,j_1,\dots,j_s\in\{1,\dots,d\},
\]
are called the \textbf{components of $T$ with respect to the chosen basis}: one plugs the dual basis elements into the first $r$ slots and the basis elements into the last $s$ slots, in all possible combinations
(``a massive array of numbers''; all indices run from $1$ to $\dim V$, the number of indices is $r+s$) \ts{3}{1}{1:33:15}\ts{3}{1}{1:33:35}\ts{3}{1}{1:33:58}\ts{3}{1}{1:34:18}\ts{3}{1}{1:34:39}\ts{3}{1}{1:35:00}\ts{3}{1}{1:35:20}\ts{3}{1}{1:35:41}\ts{3}{1}{1:36:01}\ts{3}{1}{1:36:22}.
\end{definition}
Numbers can be defined from a tensor only once a basis is chosen; ``you have to keep the bases in mind''. Why is it useful? Because, knowing the components and the bases, one can reconstruct the entire tensor \ts{3}{1}{1:36:22}\ts{3}{1}{1:36:43}\ts{3}{1}{1:37:04}.

\begin{remark}[Reconstructing a tensor from its components: the $(1,1)$ case]\label{rem:c03-comp-tensor}
Let $T$ be a $(1,1)$-tensor, with components $T^i{}_j=T(\epsilon^i,e_j)$. A map is understood if one can compute its image of any element of the domain. Expand a covector $\varphi=\sum_{i=1}^d\varphi_i\,\epsilon^i$ in the dual basis (components $\varphi_i\in\R$) and a vector $v=\sum_{j=1}^dv^j e_j$.
By multilinearity one pulls out the two sums and the real numbers:
\[
  T(\varphi,v)=T\Bigl(\sum_i\varphi_i\epsilon^i,\sum_jv^je_j\Bigr)=\sum_{i=1}^d\sum_{j=1}^d\varphi_i\,v^j\,T(\epsilon^i,e_j)=\sum_{i,j}\varphi_i\,v^j\,T^i{}_j.
\]
(The $\epsilon^i$ and $e_j$ cannot be pulled out, they are not numbers, but what remains is the definition of the components.) So the numbers $T^i{}_j$ and the bases give back the whole map
\ts{3}{1}{1:37:04}\ts{3}{1}{1:37:26}\ts{3}{1}{1:37:47}\ts{3}{1}{1:38:09}\ts{3}{1}{1:38:29}\ts{3}{1}{1:38:50}\ts{3}{1}{1:39:10}.
\end{remark}

\begin{definition}[New concepts: Einstein summation convention]\label{def:c03-einstein}
Label basis vectors with an index \emph{down} ($e_j$), basis covectors with an index \emph{up} ($\epsilon^i$), and accordingly components of vectors with indices up ($v^j$) and components of covectors with indices down ($\varphi_i$). Then indices occurring twice in a product always occur as one up and one down, and the convention is:
\emph{an index appearing once up and once down is summed over (from $1$ to $\dim V$), the sum sign being invisible}. For example
\[
  \varphi_i\,v^j\,T^i{}_j:=\sum_{i=1}^d\sum_{j=1}^d\varphi_i\,v^j\,T^i{}_j,
\]
with two invisible sum signs \ts{3}{1}{1:39:31}\ts{3}{1}{1:39:52}\ts{3}{1}{1:40:12}\ts{3}{1}{1:40:34}\ts{3}{1}{1:40:55}\ts{3}{1}{1:41:15}\ts{3}{1}{1:42:00}.
\end{definition}
\begin{aside}[The ghost sum]
Making sum signs invisible would be a bad idea, ``because if they are invisible they can be pulled out without being seen, it is like a ghost''; it only works because one always deals with multilinear maps: whether the ghost sits inside or outside the map does not matter, since for a linear map
one can pull the sum in or out. ``It saves a lot of writing effort and makes the formula appear much clearer'' \ts{3}{1}{1:41:38}\ts{3}{1}{1:42:00}.
\end{aside}
The lecture ends here: next lecture in half an hour, tutorials this afternoon \ts{3}{1}{1:42:22}.

% ==================================================================
\section{Tutorial 3}\label{sec:t03-1}

Tutorial 3, 0:57:19, \seg{3}{2}{video}, by a tutor, not Schuller. The sheet (Sheet 3) is paraphrased, not copied; see \url{https://tales.mbivert.com/on-heraeus-winter-school-tutorials/} (and \url{https://tales.mbivert.com/heraeus/Sheet_03.pdf}).
As always it starts with fifteen true/false statements, now on multilinear algebra; the grouping in three blocks of five is reconstructed from the video [?] \ts{3}{2}{0:00}\ts{3}{2}{0:24}.

\subsection*{Exercise 1: true or false}

\paragraph{(A) Vector spaces.}
\begin{itemize}
  \item Commutativity of multiplication is a vector space axiom: \textbf{false}; commutativity of the addition is, the one of multiplication follows from the underlying field \ts{3}{2}{0:45}\ts{3}{2}{1:06}.
  \item Every vector space is a matrix with only one column: \textbf{false}; the polynomials of degree at most some $N$ are a vector space and not a column \ts{3}{2}{1:06}\ts{3}{2}{1:27}.
  \item Every linear map between vector spaces can be represented by a unique square matrix: \textbf{false}; this would hold if the two dimensions were equal, but the statement says nothing about dimensions \ts{3}{2}{1:27}\ts{3}{2}{1:48}.
  \item Every vector space has a dual vector space: \textbf{true}; the linear maps from $V$ to the underlying field can always be constructed and collected \ts{3}{2}{1:48}\ts{3}{2}{2:08}.
  \item The set of everywhere positive functions on $\R$ with pointwise addition and s-multiplication is a vector space: \textbf{false}; there is no neutral element (the zero function is not positive) \ts{3}{2}{2:08}\ts{3}{2}{2:29}.
\end{itemize}

\paragraph{(B) Tensors and their components.}
\begin{itemize}
  \item The tensor product of two tensors is a tensor: \textbf{true}, ``it is defined in that way'' (the tensor product is not treated in the lecture) \ts{3}{2}{2:29}\ts{3}{2}{2:49}.
  \item One can reconstruct a tensor from its components and the corresponding basis: \textbf{true}, e.g.\ for a $(0,2)$-tensor $T=T_{ij}\,\epsilon^i\otimes\epsilon^j$ (the tutor writes $e$ for the dual basis here) \ts{3}{2}{2:49}\ts{3}{2}{3:10}.
  \item The number of indices of the tensor components depends on the dimension: \textbf{false}; the \emph{range} of the indices depends on the dimension, the number of indices ($r+s$) does not \ts{3}{2}{3:10}\ts{3}{2}{3:30}.
  \item The Einstein summation convention does not apply to tensor components: \textbf{false}; an index up and the same down hides a sum up to the dimension \ts{3}{2}{3:30}\ts{3}{2}{3:51}.
  \item A change of basis does not change the tensor components: \textbf{false}; the components are defined by plugging in the basis vectors \ts{3}{2}{3:51}\ts{3}{2}{4:13}.
\end{itemize}

\paragraph{(C) Bases.}
\begin{itemize}
  \item Given a basis of a $d$-dimensional $V$, one finds exactly $d^2$ dual bases of $V^*$: \textbf{false}; there is precisely one dual basis, given by $\epsilon^a(e_b)=\delta^a{}_b$, a linear system with exactly one solution \ts{3}{2}{4:13}\ts{3}{2}{4:34}\ts{3}{2}{4:56}\ts{3}{2}{5:16}.
  \item Removing one basis vector of a basis of $V$ gives a basis of a $(d-1)$-dimensional vector space: \textbf{true}; e.g.\ removing a vector of the standard basis of $\R^3$ leaves a basis of a plane (the span of the others, a subspace) \ts{3}{2}{5:16}\ts{3}{2}{5:37}.
  \item The continuity of a map $f\colon V\to W$ depends on the choice of a basis of $W$: \textbf{false}; continuity is defined for topological spaces, not through a basis \ts{3}{2}{5:37}\ts{3}{2}{5:58}.
  \item Given a basis of $V$ one can construct the components of the elements of $V^*$: \textbf{true}; the components are defined by plugging in the basis \ts{3}{2}{5:58}\ts{3}{2}{6:21}.
  \item Given a basis, each vector of $V$ can be reconstructed from its components: \textbf{true}, as seen in the lecture, vectors being $(1,0)$-tensors \ts{3}{2}{6:21}\ts{3}{2}{6:41}\ts{3}{2}{7:02}.
\end{itemize}

\subsection*{Exercise 2: $\R^3$, a linear map, a covector and polynomials}

We build $V=\R^3$ as a vector space and relate it to the polynomials of lecture 3 (here of degree at most $2$) \ts{3}{2}{7:02}\ts{3}{2}{7:24}.
\begin{enumerate}
  \item \emph{Vector space.} Equip $\R^3$ with entrywise addition $(a,b,c)\oplus(d,e,f):=(a+d,b+e,c+f)$ and s-multiplication $\lambda\odot(a,b,c):=(\lambda a,\lambda b,\lambda c)$ (circled symbols, to distinguish them from the operations in the entries). One checks the eight axioms;
        each follows from the corresponding property of $\R$ in each entry (neutral element $(0,0,0)$, inverse $(-a,-b,-c)$, unit $1$). The full check, axiom by axiom, is moved to the appendix (appendix~
ef{app:exam} (section ``Chapter 3'')) \ts{3}{2}{7:24}\ts{3}{2}{18:47}.
  \item \emph{A linear map.} $D\colon V\to V$, $(a,b,c)\mapsto(b,2c,0)$. It is linear: $D(\lambda v\oplus w)=\lambda\odot D(v)\oplus D(w)$, computed entrywise (the entries of $D(\lambda(a,b,c)+(d,e,f))$ are $\lambda b+e$, $2(\lambda c+f)$, $0$, which is $\lambda D(a,b,c)+D(d,e,f)$); the tutor keeps a tilde on the arrow,
        $D\colon V\xrightarrow{\sim}V$. Likewise $D\circ D\colon(a,b,c)\mapsto(2c,0,0)$ is linear \ts{3}{2}{19:30}\ts{3}{2}{19:51}\ts{3}{2}{20:13}\ts{3}{2}{20:34}\ts{3}{2}{20:54}\ts{3}{2}{21:15}\ts{3}{2}{21:36}\ts{3}{2}{21:57}\ts{3}{2}{22:18}\ts{3}{2}{23:08}\ts{3}{2}{23:28}\ts{3}{2}{23:48}\ts{3}{2}{24:09}\ts{3}{2}{24:29}\ts{3}{2}{24:50}\ts{3}{2}{25:10}\ts{3}{2}{25:30}.
  \item \emph{A covector.} $i\colon V\to\R$, $(a,b,c)\mapsto a+\tfrac12b+\tfrac13c$. It is linear (the factor $\lambda$ floats out of $a+\tfrac12b+\tfrac13c$), so $i\in V^*$ since $V^*$ is the set of linear maps $V\to\R$ \ts{3}{2}{25:51}\ts{3}{2}{26:12}\ts{3}{2}{26:35}\ts{3}{2}{26:55}\ts{3}{2}{27:16}\ts{3}{2}{27:38}\ts{3}{2}{28:00}.
  \item \emph{A $(0,2)$-tensor.} $g\colon V\times V\to\R$ is defined on the sheet [formula not read out in the video; by the link below it is the inner product of the polynomials in the coordinates, ?]. One checks linearity in each entry, e.g.\ $g(\lambda v+w,x)=\lambda g(v,x)+g(w,x)$, and the same for the second entry; hence $g$ is multilinear,
        a $(0,2)$-tensor \ts{3}{2}{28:23}\ts{3}{2}{28:43}\ts{3}{2}{29:05}\ts{3}{2}{29:26}\ts{3}{2}{29:49}.
  \item \emph{Why these maps?} Let $P_2$ be the quadratic polynomials $p(x)=a+bx+cx^2$ and define $J\colon P_2\to\R^3$, $p\mapsto(a,b,c)$ with inverse $J^{-1}(a,b,c)=(x\mapsto a+bx+cx^2)$, a linear bijection between the vector spaces. The differentiation operator $\Delta$ of the lecture gives $\Delta p=b+2cx+0x^2$, which is $D$ moved by $J$:
        \[
          D=J\circ\Delta\circ J^{-1}\qquad\text{(the diagram: $P_2\xrightarrow{\Delta}P_2$ above, $\R^3\xrightarrow{D}\R^3$ below, $J$ the vertical arrows).}
        \]
        Likewise the integral $I(p)=\int_0^1p\,\mathrm dx=a+\tfrac12b+\tfrac13c$ of the lecture becomes $i=I\circ J^{-1}$ (the tutor writes ``$I=\dots J$'', garbled [?]); and the same for $g$ and the inner product of the lecture. So the structure-preserving map $J$ connects the two vector spaces and one map can be written in terms of the other
        \ts{3}{2}{30:09}\ts{3}{2}{30:29}\ts{3}{2}{30:49}\ts{3}{2}{31:10}\ts{3}{2}{31:32}\ts{3}{2}{31:53}\ts{3}{2}{32:14}\ts{3}{2}{32:34}\ts{3}{2}{32:55}\ts{3}{2}{33:16}\ts{3}{2}{33:40}\ts{3}{2}{34:00}\ts{3}{2}{34:21}\ts{3}{2}{34:42}\ts{3}{2}{35:04}\ts{3}{2}{35:25}\ts{3}{2}{35:46}\ts{3}{2}{36:07}\ts{3}{2}{36:28}\ts{3}{2}{36:48}.
\end{enumerate}
(``$J^{-1}$'' is called $J$ with an exponent $-1$ in the transcript and one of the ``$J$'' arrows is garbled; the reconstruction above follows the diagram described by the tutor \ts{3}{2}{34:42}\ts{3}{2}{35:04}.)

\subsection*{Exercise 3: indices, symmetrisation, antisymmetrisation}

Let $V$ be $d$-dimensional with a basis $e_1,\dots,e_d$ and dual basis $\epsilon^1,\dots,\epsilon^d$, $A$ a $(2,0)$-tensor and $B$ a $(0,2)$-tensor \ts{3}{2}{36:48}\ts{3}{2}{37:11}.
\begin{enumerate}
  \item \emph{Components:} $A^{ab}:=A(\epsilon^a,\epsilon^b)$ and $B_{ab}:=B(e_a,e_b)$ (as in \cref{def:c03-comp-tensor}) \ts{3}{2}{37:33}\ts{3}{2}{37:54}\ts{3}{2}{38:14}.
  \item \emph{Einstein convention in one's own words:} indices that appear pairwise, once up and once down, are summed over; e.g.\ in a product of tensors $T$ and $S$ with a repeated index up and down one sums over it \ts{3}{2}{38:14}\ts{3}{2}{38:36}\ts{3}{2}{38:59}\ts{3}{2}{39:20}\ts{3}{2}{39:41}.
  \item \emph{Antisymmetrisation.} Define $B_{[ab]}:=\tfrac12(B_{ab}-B_{ba})$. (i) Exchanging the two indices gives a minus sign, $B_{[ba]}=\tfrac12(B_{ba}-B_{ab})=-B_{[ab]}$. (ii) The brackets can be moved from one index pair to the other:
        \[
          A^{[ab]}B_{ab}=\tfrac12\bigl(A^{ab}B_{ab}-A^{ba}B_{ab}\bigr)=\tfrac12\bigl(A^{ab}B_{ab}-A^{ab}B_{ba}\bigr)=A^{ab}B_{[ab]},
        \]
        where in the second term the (summed, hence arbitrary) dummy indices are renamed $a\leftrightarrow b$; ``it does not matter whether they are called $\mu$, $\nu$ or whatever''. This saves a lot of calculation later \ts{3}{2}{40:02}\ts{3}{2}{40:23}\ts{3}{2}{40:44}\ts{3}{2}{41:04}\ts{3}{2}{41:25}\ts{3}{2}{41:46}\ts{3}{2}{42:07}\ts{3}{2}{42:29}\ts{3}{2}{42:50}\ts{3}{2}{43:12}.
  \item \emph{Symmetrisation.} Define $B_{(ab)}:=\tfrac12(B_{ab}+B_{ba})$. (i) $B_{(ab)}=B_{(ba)}$ since addition of reals is commutative. (ii) Likewise $A^{(ab)}B_{ab}=A^{ab}B_{(ab)}$ (same renaming of dummy indices) \ts{3}{2}{43:32}\ts{3}{2}{43:53}\ts{3}{2}{44:14}\ts{3}{2}{44:36}\ts{3}{2}{44:57}\ts{3}{2}{45:18}\ts{3}{2}{45:38}.
  \item \emph{A useful result:} if one index pair is summed with an antisymmetric and a symmetrised object the result is zero:
        \[
          A^{[ab]}B_{(ab)}=\tfrac12\bigl(A^{ab}B_{(ab)}-A^{ba}B_{(ab)}\bigr)=\tfrac12\bigl(A^{ab}B_{(ab)}-A^{ab}B_{(ba)}\bigr)=0,
        \]
        since $B_{(ab)}=B_{(ba)}$ (rename the dummy indices in the second term, use (i) of the symmetrisation) \ts{3}{2}{45:38}\ts{3}{2}{45:58}\ts{3}{2}{46:19}\ts{3}{2}{46:40}\ts{3}{2}{47:00}\ts{3}{2}{47:23}\ts{3}{2}{47:45}\ts{3}{2}{48:07}.
\end{enumerate}
The tutor names the square brackets \emph{antisymmetrisation} and the round brackets \emph{symmetrisation} of index pairs \ts{3}{2}{48:07}. (The tutor writes the brackets moved to the upper indices; the displayed chain is a reconstruction of the same steps, the original line is garbled [?].)

\subsection*{Exercise 4: linear maps as $(1,1)$-tensors}

This is the construction of the student's question in the lecture (\cref{sec:c03-tensors}). Let $\phi\colon V^*\to V^*$ be linear \ts{3}{2}{48:30}\ts{3}{2}{48:50}.
\begin{enumerate}
  \item From $\phi$ construct a $(1,1)$-tensor $T_\phi\colon V^*\times V\to\R$: take $\sigma\in V^*$, $v\in V$, apply $\phi$ to $\sigma$ (again an element of $V^*$, i.e.\ a map $V\to\R$) and feed it $v$:
        $T_\phi(\sigma,v):=\phi(\sigma)(v)$ \ts{3}{2}{49:11}\ts{3}{2}{49:31}\ts{3}{2}{49:52}\ts{3}{2}{50:13}.
  \item The other way round, from a $(1,1)$-tensor $T$ construct a linear map $\phi_T\colon V^*\to V^*$: for $\sigma\in V^*$ put $\phi_T(\sigma):=T(\sigma,\,\cdot\,)$ with an open slot; plugging in a vector $v$ gives a real number, so this is an element of $V^*$ as it should be
        \ts{3}{2}{50:13}\ts{3}{2}{50:33}\ts{3}{2}{50:55}\ts{3}{2}{51:15}\ts{3}{2}{51:36}.
  \item The definitions are good: $T_{\phi_T}=T$, because $T_{\phi_T}(\sigma,v)=\phi_T(\sigma)(v)=T(\sigma,\,\cdot\,)(v)=T(\sigma,v)$; and $\phi_{T_\phi}=\phi$, because $\phi_{T_\phi}(\sigma)=T_\phi(\sigma,\,\cdot\,)=\phi(\sigma)$ (no $v$ left in the expression)
        \ts{3}{2}{51:36}\ts{3}{2}{51:57}\ts{3}{2}{52:18}\ts{3}{2}{52:39}\ts{3}{2}{53:01}\ts{3}{2}{53:21}\ts{3}{2}{53:42}\ts{3}{2}{54:03}\ts{3}{2}{54:23}\ts{3}{2}{54:44}\ts{3}{2}{55:05}.
  \item Conclusion: one can \emph{consider} linear maps $V^*\to V^*$ as $(1,1)$-tensors (``consider'' is better than ``is'', since it is not precisely the same, but the one can be reconstructed from the other and the correspondence is well defined) \ts{3}{2}{55:05}\ts{3}{2}{55:26}\ts{3}{2}{55:47}\ts{3}{2}{56:09}.
\end{enumerate}
A last note: for a map $\psi\colon V\to V$ one can also think of a $(1,1)$-tensor with a similar construction, but this is more sophisticated, because one has to use that $V^{**}$ is isomorphic to $V$, which would take about a page and is left to the audience \ts{3}{2}{56:30}\ts{3}{2}{56:51}\ts{3}{2}{57:11}.

% ==================================================================
\section*{Key concepts}
\begin{keyconcepts}
\textbf{Retrieval questions.}
\begin{enumerate}
  \item Why does spacetime not carry a vector space structure, and why study vector spaces at all?
  \item State the eight axioms of a real vector space. Which plus and which times is which in the axioms with $\lambda,\mu$?
  \item Is the polynomial $x^2$ a vector? What does the question mean?
  \item When is a map linear? Why can it be called linear only after the structures on domain and target are given?
  \item What is $\operatorname{Hom}(V,W)$ and how is it made into a vector space? What is $V^*$ and a covector? Give an example of a covector on $P$.
  \item Define an $(r,s)$-tensor. Why are vectors $(1,0)$-tensors, covectors $(0,1)$-tensors and linear maps $(1,1)$-tensors? Where is $\dim V<\infty$ used?
  \item Why should one avoid choosing a basis? What is the cost, and how does it relate to charts?
  \item What is a Hamel basis, a Schauder basis, the dimension? Write down the dual basis of $\{1,x,x^2,x^3\}$.
  \item Define the components $T^{i_1\dots i_r}{}_{j_1\dots j_s}$ and reconstruct a $(1,1)$-tensor from them. State the Einstein convention.
\end{enumerate}
\textbf{Concepts.}
\begin{itemize}
  \item Vector space $(V,+,\cdot)$ with eight axioms (\cref{def:c03-vs}); linear map $f\colon V\xrightarrow{\sim}W$ (\cref{def:c03-linear}); $\operatorname{Hom}(V,W)$ (\cref{def:c03-hom}).
  \item Dual space $V^*=\operatorname{Hom}(V,\R)$ and covectors (\cref{def:c03-dual}); example $I(p)=\int_0^1p$ (\cref{ex:c03-int}); the gradient is a covector.
  \item $(r,s)$-tensor as multilinear map $(V^*)^r\times V^s\to\R$ (\cref{def:c03-tensor}); vectors, covectors, linear maps, inner products as tensors (\cref{thm:c03-vec-tensors}, \cref{ex:c03-ip}).
  \item Hamel basis, dimension (\cref{def:c03-basis,def:c03-dim}); components of vectors (\cref{rem:c03-components}); dual basis $\epsilon^a(e_b)=\delta^a{}_b$ (\cref{def:c03-dualbasis}).
  \item Components of tensors (\cref{def:c03-comp-tensor,rem:c03-comp-tensor}) and the Einstein summation convention (\cref{def:c03-einstein}).
  \item (Anti)symmetrisation of indices $A^{[ab]}B_{(ab)}=0$ (\cref{sec:t03-1}).
\end{itemize}
\end{keyconcepts}

\flushside
\end{paracol}
